\section{Scale estimates for general local arrays}
\label{sec:scale}

We first state the scale estimates that will be used to choose the
temperature.  These are results about local random potentials, rather
than about a particular one-edge distribution.  We use the exact energy
charge, local transformation, moment estimates, and variance coordinate
of \cite[Sections~2--5]{R}.  The hypotheses below explain why those
results apply to the present law \(\rho\), although the main theorem of
\cite{R} concerns symmetric two-valued disorder.  The pure-model inputs
behind that construction are developed in \cite[Sections~2--3]{S}.

\subsection{The reference field and its energy charge}

Let \(\E_0\) denote expectation under the full-plane critical Ising
law at \(K_0\), together with a locally finite family of independent,
finite-valued reference variables attached to specified locations.
These variables are part
of the Gibbs field: they are integrated or conditioned along with the
spins.  Spin reversal leaves them unchanged.  Their product reference
law, including their locations and labels, respects lattice translations
and square symmetries.  At the initial scale there are no such variables.

For an edge \(e=\{x,y\}\), put
\[
 b_e=\sigma_x\sigma_y-\E_0(\sigma_x\sigma_y),
 \qquad B_x=\frac14\sum_{e\ni x}b_e.
\]
The theorem ``Exact pure energy charge'' in \cite[Section~2]{R}
provides the linear functional
\begin{equation}
 Q(f)=\kappa^{-1}\lim_{|x|\to\infty}|x|^2\E_0(fB_x),
 \qquad \kappa=\pi^{-2},
\label{eq:energy-charge}
\end{equation}
on bounded, spin-flip-even local observables.  The reference variables
are included in this expectation.  Thus \(Q\) measures the energy
component of a local perturbation through its correlation with a
distant energy observable.  The functional kills constants,
is invariant under lattice translations and square symmetries, and
satisfies \(Q(b_e)=1\).  Pure conditional integration across a surrounding
square preserves \(Q\) exactly.  If the support of \(f\) lies in a square
of radius \(r\), then
\begin{equation}
 |Q(f)|\le C r\|f\|_\infty.
 \label{eq:charge-bound}
\end{equation}
This normalization permits charges at successive scales to be compared
without accumulating a normalization error.

\subsection{Arrays and the exact scale transformation}

At an integer spacing \(s\), a local array is a family
\(f=(f_i)_{i\in\Z^2}\) of random functions based at the cell centers
\(si\).  We impose the following hypotheses.  For every environment,
each \(f_i\) is real, bounded, spin-flip even, and pure centered:
\(\E_0 f_i=0\).  Both its field support and the primitive environmental
variables determining it lie within distance \(Cs\) of \(si\), in
supremum norm, for the fixed support constant of the scale construction.
The primitive environmental variables are independent.  The law of the
array is invariant under translations of the \(s\)-grid and the
symmetries of a square, acting also on reference-variable locations and
labels.  These are symmetries in distribution; an individual array need
not have them.  Finally, \(\E\|f_0\|_\infty^{40}<\infty\).

Here and throughout the scale argument, \(\E\) averages the bond
environment and the independent priorities and thresholds used to
construct the transformation.  These choices are held fixed when
sampling the Gibbs field.  They are distinct from the reference variables
averaged by \(\E_0\).  In a finite volume, the array adds
\(\sum_i f_i\) to the logarithm of the reference density, with all old
exterior field variables prescribed.  The sum includes the potentials
meeting the volume.

The proposition ``The local scale map'' in \cite[Section~3]{R} gives,
for a sufficiently large odd integer \(L\) and a threshold parameter
\(b>0\), a map from such arrays at spacing \(s\) to arrays at spacing
\(s'=Ls\) with the same support constant.  The construction uses
independent continuous priorities and independent environmental
thresholds uniform on \([b,2b]\) to determine its local choices and to
flag large input terms.  The map introduces normalized conditional
likelihoods for fresh reference coins, then uses conditional integrations
in separated squares to replace local logarithmic interactions.  These
are changes of the interaction functions: the enlarged field alphabet
retains every old spin and reference variable.
Centering the resulting potentials changes only a scalar factor.  For
each fixed environment and arbitrary old exterior assignment, these
operations preserve the exterior-integrated weight exactly.  The same
formulas and full-plane centering constants apply in embedded torus
charts.  In particular, the construction is compatible with a flat spin
twist: a change of trivialization in a local chart is a constant spin
flip, under which all potentials are invariant.

We will use three quantitative and local features of that proposition,
the corollary ``Local implementation'', and the proposition ``A coupling
of consecutive fields'' in the same section.  There are
fixed numerical constants \(K,K_1\) and a constant \(C_L\), depending
on the fixed step parameters, such that
\begin{equation}
 \|f'_j\|_\infty\le C_L
 \sum_{i:\,\|si-s'j\|_\infty\le Ks'}\|f_i\|_\infty.
 \label{eq:raw-scale-bound}
\end{equation}
In a finite free chart with prescribed old exterior field data, the
transformation agrees with the full output specification after erosion
by \(K_1s'\).  This assertion includes annular charts.  Every fresh
reference variable with a nonconstant likelihood is sampled as part of
the construction.  Exactness does not assert the same input kernel if
fresh coins are independently prescribed before the construction;
disintegration over the sampled output exterior remains available.
Finally, consecutive fields can be
coupled to agree outside squares of radius \(5s'\) about marked parent
centers.  Conditional on the enlarged environment, these marks are
independent of one another and of the output Gibbs sample, with local
parameters satisfying
\begin{equation}
 p_j\le \min\left\{1,C_L
 \sum_{i:\,\|si-s'j\|_\infty\le Ks'}\|f_i\|_\infty\right\},
 \qquad \E p_j\le C_L(\E\|f_0\|_\infty^{40})^{1/40}.
 \label{eq:scale-marks}
\end{equation}
Both the coupling and these bounds are uniform over the prescribed old
exterior assignments.  These exactness and locality statements will
later connect the moment estimates to crossing events.

\subsection{Moment and variance coordinates}

Write \(q_s=Q/s\), and define
\begin{equation}
 \begin{aligned}
 X&=(\E\|f_0\|_\infty^{40})^{1/40},
 &Y&=\|\E f_0\|_\infty,\\
 m&=\E q_s(f_0),
 &W&=\sum_{i\in\Z^2}\Cov\bigl(q_s(f_0),q_s(f_i)\bigr).
 \end{aligned}
 \label{eq:scale-coordinates}
\end{equation}
All covariances here are environmental.  The covariance sum is finite
by the dependence hypothesis.  It is nonnegative because, for growing
squares \(\Lambda\subset\Z^2\),
\[
 W=\lim_{\Lambda\uparrow\Z^2}\frac1{|\Lambda|}
       \Var\left(\sum_{i\in\Lambda}q_s(f_i)\right).
\]
The mean charge \(m\) is the expanding coordinate of the map.  The
variance density \(W\) is marginal at its linearization.  To see its
decay, one needs a coordinate that also retains local cubic and quartic
terms.

The construction in \cite[Section~5, Definition ``Variance coordinate'']{R}
begins with a finite random
insertion \(F=\sum_{i\in I}t_i f_i\), where \(I\subset\Z^2\) is finite
and the \(t_i\) are real parameters.  Its normalized charge is
\[
 T(F)=q_s\left(\frac{e^F}{\E_0e^F}\right).
\]
We take Taylor coefficients at zero for each fixed environment; the
expansions below use only finite moments, not environmental exponential
integrability.
Since \(\E_0F=0\) and \(q_s\) kills constants, the first three
homogeneous terms are
\[
 T_1=q_s(F),\qquad T_2=\tfrac12q_s(F^2),\qquad
 T_3=\tfrac16q_s(F^3)-\tfrac12q_s(F)\E_0F^2.
\]
The terms of degrees two, three, and four in the formal environmental variance
of \(T(F)\) are therefore
\[
 \Var(T_1),\qquad 2\Cov(T_1,T_2),\qquad
 \Var(T_2)+2\Cov(T_1,T_3).
\]
The cubic term need not vanish for a nonsymmetric law.  We retain it
while estimating the degree-four drift: an uncontrolled cubic error
would be larger than the quartic change we seek to determine.

The coordinate \(v_H\) is this variance polynomial through degree four, per cell,
with a cutoff on the diameter of the input labels.  More precisely,
choose a dependence range \(R_0\) in cell units so that sets separated
by more than \(R_0\) use disjoint environmental variables and have
disjoint declared supports.  For an ordered tuple
\(\boldsymbol i=(i_1,\ldots,i_d)\), \(2\le d\le4\), give each
occurrence its own indeterminate, including repeated indices, and set
\begin{equation}
 C_d(\boldsymbol i)=
 \left.\partial_{t_1}\cdots\partial_{t_d}\,
 \Var\left[
 q_s\left(
 \frac{\exp(\sum_{a=1}^d t_a f_{i_a})}
      {\E_0\exp(\sum_{a=1}^d t_a f_{i_a})}
 \right)\right]\right|_{\boldsymbol t=0}.
 \label{eq:variance-coefficients}
\end{equation}
The derivatives mean coefficient extraction from the displayed variance
polynomial.  For \(H>R_0\), define
\begin{equation}
 v_H=\sum_{d=2}^4\frac1{d!}
 \sum_{i_2,\ldots,i_d\in\Z^2}
 \1_{\{\diam(0,i_2,\ldots,i_d)\le H\}}
 C_d(0,i_2,\ldots,i_d),
 \label{eq:variance-coordinate}
\end{equation}
where diameter is Euclidean diameter in cell units.  This is a real
polynomial coordinate, not necessarily positive on an arbitrary array.
Its quadratic part is exactly \(W\).

\begin{proposition}[Imported scale estimates]
\label{prop:scale-estimates}
For the local arrays just specified, the scale-step parameters can be
chosen in the order \(L\) sufficiently large, then \(b>0\)
sufficiently small, then \(x_*>0\).  There are constants
\(0<\lambda<1\) and \(C_1<\infty\), independent of the spacing, such
that every input with \(X\le x_*\) satisfies
\begin{equation}
 \begin{split}
 |m'-Lm|&\le C_1X^2,\\
 Y'&\le\lambda Y+C_1(|m|+X^2),\\
 X'&\le\lambda X+C_1(\sqrt W+Y).
 \end{split}
 \label{eq:one-step-estimates}
\end{equation}
Primes denote the output coordinates, with charge normalization
\(q_{Ls}\).  For every fixed \(H>R_0\),
\begin{equation}
 |v_H-W|\le C_HX^3\qquad (X\le1).
 \label{eq:variance-near-W}
\end{equation}
Furthermore, fix finite constants \(A,B\).  For every sufficiently
large fixed \(H\), whenever \(0<z\le1\), \(X\le Az\), and
\(Y\le Bz^2\), one has
\begin{equation}
 v'_H-v_H=-2\kappa W^2
       \sum_{y\in\Z^2:\,H<|y|\le LH}|y|^{-2}
       +e_H z^4+O_H(z^5).
 \label{eq:marginal-variance-flow}
\end{equation}
Here \(|e_H|\) is bounded by a deterministic quantity tending to zero
as \(H\to\infty\).  The constants may depend on \(A,B\) and the fixed
step parameters, but are uniform in the lattice spacing.
\end{proposition}

These are respectively the propositions ``One-step estimates'' and
``Marginal variance flow'' of \cite[Sections~4--5]{R}, with the
near-\(W\) estimate included in the latter.  Their hypotheses are the
array hypotheses above.  The required symmetry is spatial symmetry in
law, together with spin-flip evenness of the potentials; there is no
assumption of invariance under negating the bond randomness.  In the
variance calculation, independence eliminates isolated centered
environmental occurrences, while the cubic terms in
\eqref{eq:variance-coordinate} are kept.  Thus an asymmetric law does
not require any alteration of the imported estimates.  The next section
verifies their hypotheses for both temperature families and uses the
negative term in \eqref{eq:marginal-variance-flow} to keep the entire
iteration small.
